% Ad Quintum Fratrem 3.3 (ad-quintum-fratrem-03-03) --- English translation
% Translated by Claude (Anthropic) from the Latin (Perseus).
% Style guide: STYLE.md.

\ciceroLetterOpener{MARCVS QUINTO FRATRI SALUTEM.}

\ciceroSection{1} Let my secretary's hand be a sign of my occupations. There is, you may know, no day on which I do not speak for a defendant. So whatever I get done or contemplate, I throw it for the most part into walking-time. Our affairs stand thus, and the domestic ones as we wish. The boys are well, they learn enthusiastically, they are taught diligently, they love us and one another. The polishing of both our houses is in hand; but yours, the country business at Arcanum and Laterium, is now done. Besides, on the water and the road I have left nothing out in a certain letter, but kernel by kernel wrote it through to you. But that worry vexes and troubles me extremely, that for an interval now of more than fifty days nothing from you, nothing from Caesar, nothing out of those parts has flowed in --- no letter, not even rumour. And the sea over there and the land are now troubling me, and I do not cease (as happens in affection) to imagine the things I least wish to. So I do not now ask you to write to me about yourself, about those affairs (for you never miss the chance when you can write); but this I want you to know, that I have scarcely ever waited for anything as, when I was writing this, I have been waiting for your letter.

\ciceroSection{2} Now learn what is going on in public life. The days of the elections are being knocked off one by one through formal announcements of bad omens, with the great willingness of all good men --- so much hatred do the consuls labour under for the suspicion of bargains promised by the candidates. The four consular candidates are all on trial. The cases are difficult, but we will strain ourselves to keep our Messala safe, which is linked also to the safety of the rest. Gabinius has been put on trial for canvassing by P. Sulla, with his stepson Memmius, his brother Caecilius, and his son Sulla as supporting counsel. L. Torquatus argued against, and, though everyone was in his favour, did not prevail.

\ciceroSection{3} You ask what will become of Gabinius. We shall know about the \textit{maiestas} charge within three days; in that trial he is pressed by hatred of every kind, is wounded most by the witnesses, and is using the most lukewarm prosecutors; the panel is a mixed one, the president of the court, Alfius, is weighty and firm, Pompey is vehement in his lobbying of the jurors. What is going to happen I do not know; a place for him in the State, however, I do not see. I am holding my spirit moderate towards his ruin, and entirely mild towards the outcome.

\ciceroSection{4} You have on the whole everything. One thing I will add: your Cicero, who is ours as well, is at the height of his enthusiasm with his rhetor Paeonius, a very practised and good man, I think. But our way of training is a touch more cultivated and \textit{[Greek: thetik\=oteron]} (more inclined to general theses), as you do not fail to know. So I do not want either Cicero's progress hindered or that schooling broken off, and the boy himself appears to be more led on and delighted by the declamatory genre. And since we ourselves were also in that, let us allow him to go by our paths; for I am confident he will arrive at the same place. Still, if at any point we shall take him out to the country with us, we will lead him into this way and habit of ours. For great is the reward you have set before us, which by our own fault we will certainly never fall short of. In what places, and with what hope, you are to winter, I should like you to write to me as carefully as possible.
